\chapter{Implementation and Evaluation of the Anomaly Detection and Fault Diagnosis Pipeline}
\label{ch:implementation}

\section{Introduction}
This chapter details the implementation of the anomaly detection and fault diagnosis methods explored during this project. It covers the model architectures, training procedure and results
obtained from the evaluation on the Costa dataset.

\section{Tools and Experimental Setup}
\subsection{Software and Libraries}
The implementation relies on a small set of libraries and services, grouped by their role in the end-to-end pipeline.

\subsubsection{ML/DL core}
\par\noindent\textbf{Python:} The primary programming language used for data processing, model development, and evaluation.
\par\noindent\textbf{PyTorch:} The deep learning framework used for implementing and training the neural models.
\par\noindent\textbf{Scikit-learn:} Used for preprocessing, classical baselines, and evaluation metrics.
\par\noindent\textbf{SciPy:} Scientific computing utilities used for numerical routines and statistical operations required by parts of the pipeline.
\par\noindent\textbf{Optuna:} Hyperparameter optimization framework used to tune model configurations under a shared HPO protocol.

\subsubsection{Data engineering and visualisation}
\par\noindent\textbf{Polars:} A high-performance DataFrame library used for efficient tabular data transformations in the preprocessing and feature engineering stages.
\par\noindent\textbf{Matplotlib:} Used to generate diagnostic plots and result figures (e.g., curves and confusion matrices).

\subsubsection{Experiment tracking and reproducibility}
\par\noindent\textbf{GitHub:} Used for version control and collaborative code management.
\par\noindent\textbf{Data Version Control (DVC):} Used to version datasets and large artifacts, and to encode the multi-stage pipeline as reproducible dependency graphs.
\par\noindent\textbf{MLflow:} Free software used for experiment tracking, metrics logging, and model versioning.
\par\noindent\textbf{DagsHub:} Used as a centralized platform for hosting the DVC remote and the MLflow tracking server, enabling traceability of runs and artifacts.
\par\noindent\textbf{Google Colab:} Used as a managed GPU-backed execution environment for deep learning training and hyperparameter search.

\subsection{Experimental Setup}\label{sec:implementation-experimental-setup}
This section summarizes the software stack and execution environment used to implement and evaluate the proposed Fault Detection and Diagnosis (FDD) pipeline. The goal is to ensure that the reported results can be reproduced from a versioned codebase, a traceable dataset state, and a documented configuration.

\paragraph{Codebase and execution contract.}
All experiments are implemented as a Python (\(\geq 3.11\)) monorepository, where each pipeline stage is exposed through a module entry point executed with a unified invocation pattern (e.g., \texttt{uv run python -m src.data.ingestion}). The \texttt{uv} build system is used as the only package and environment manager, ensuring consistent dependency resolution across local and cloud execution.

\paragraph{Configuration management.}
To minimize implicit assumptions and to enable systematic comparisons, all data and model decisions are externalized into YAML configuration files. The data configuration defines the active dataset, split policy, preprocessing options, and feature profiles, while the model configuration defines the active model lineup, hyperparameter optimization (HPO) settings, and detector-specific post-processing (e.g., anomaly threshold calibration). When multiple configuration scopes apply, parameters are resolved through a deterministic hierarchy (global defaults, dataset overrides, task directives, then profile overrides), which prevents silent overrides and simplifies run auditing.

\paragraph{Data versioning and pipeline reproducibility.}
Large artifacts (raw data, processed tables, and selected model checkpoints) are managed with Data Version Control (DVC). The end-to-end workflow is encoded as a multi-stage DVC pipeline (ingestion, splitting, preprocessing, featurization, and training), such that intermediate artifacts are regenerated from their upstream dependencies and tracked through content hashes. The DVC remote storage is hosted on DagsHub, while Git stores the lightweight \texttt{.dvc} pointers and the pipeline definition. This separation keeps the repository manageable while preserving traceability between the reported metrics and the exact data state that produced them.

\paragraph{Compute environments and execution discipline.}
Two execution environments are used depending on workload characteristics. CPU-oriented stages, including ingestion, exploratory analysis, and the classical machine learning baselines, are executed on a local machine. Deep learning training and HPO are executed on Google Colab with GPU acceleration (typically NVIDIA T4 or A100). In Colab sessions, a persistent Google Drive workspace is mounted and used as an artifacts root directory, so that Optuna studies, checkpoints, metrics, and plots survive runtime disconnections. This split avoids wasting GPU resources on CPU-bound data manipulation and focuses accelerated compute on the training loops that benefit from it.

\paragraph{Hyperparameter optimization protocol.}
% Model selection is conducted with Optuna using the Tree-structured Parzen Estimator (TPE) sampler. Classical models are tuned with fixed trial budgets, where each trial is fast and therefore pruning is not required. Deep learning models use the same Optuna interface but can additionally enable early stopping and pruning strategies, because training a single trial can span multiple epochs. For all tasks, the validation partition is the only split used for HPO decisions (metric optimization and threshold selection), while the test partition is reserved for final reporting.
Each model evaluated in this work exposes a set of parameters that govern its capacity and training dynamics (learning rate, regularization coefficients, architectural dimensions) but that training itself cannot determine. These must be selected through a separate optimization process driven by a held-out validation signal. A systematic, automated search is necessary: the number of candidate configurations is too large for manual exploration, and deep learning models in particular carry non-trivial trial costs that penalize exhaustive strategies.

Two complementary optimization strategies are employed depending on the model family. For most classical and deep learning models, optimization is performed using Optuna \citep{akiba2019optuna} with the Tree-structured Parzen Estimator (TPE) sampler. Rather than enumerating a predefined grid of configurations as grid search does \citep{bergstra2012random}, or evolving a population of candidates as genetic algorithms require, TPE builds a probabilistic model of the objective from the history of completed trials and focuses subsequent sampling on the regions associated with good outcomes. This sequential, model-guided search reaches competitive configurations in far fewer trials than either exhaustive or population-based alternatives, which is decisive when individual trials involve full training runs of deep architectures.

The second strategy is the Good-point-set Vectorised Snow Ablation Optimiser (GVSAO), a population-based metaheuristic derived from the snow ablation algorithm \citep{GTBADGVSAOTransformerBiLSTMbasedTimeseries2026}. GVSAO initialises a population of candidate solutions using a good-point-set method that ensures uniform coverage of the search space from the first generation. Each candidate is evaluated on the objective, and the population is refined over successive generations through dual-population dynamics that balance exploration and exploitation: one sub-population performs local refinement around high-quality solutions while the other maintains diversity by searching less-visited regions. Adaptive melting factors modulate the intensity of this competition across generations, and periodic oscillation mutations inject controlled randomness to prevent premature convergence. Unlike TPE, which builds a probabilistic model of the objective surface from trial history, GVSAO directly navigates the parameter space through coordinated agent movement, making it suitable when the evaluation budget is small and the overhead of maintaining a surrogate model would outweigh its benefit.

Each HPO study runs on the fixed original train/val split and completes before the K-fold cross-validation begins, so the winning configuration is fixed for all folds. The test partition is invisible throughout. The optimization objective is the validation macro per-class PR-AUC for anomaly detection and the validation macro F1 for fault classification, consistent with the headline metrics of Section~\ref{sec:evaluation-methodology}. Among configurations whose scores fall within one standard deviation of each other, the simpler configuration is preferred as a tiebreaker.

The pruning strategy is adapted to trial cost. For classical ML models, no pruner is applied: trials are inexpensive and no intermediate epoch-level signal is available. For deep learning models, a Median pruner is added: a trial is terminated when its intermediate validation score falls below the median of all completed trials at the same epoch, after a warmup period during which pruning is suspended. Within-trial early stopping patience scales with the trial epoch budget, preventing premature termination on short trial budgets.


\paragraph{Experiment tracking and artifacts.}
All runs are tracked with MLflow and logged to a centralized DagsHub tracking server. Each run records the active dataset and split identifier, the selected feature profile, the final HPO parameters, and the computed validation and test metrics. In addition, lightweight artifacts are stored alongside metrics, including serialized preprocessing objects (e.g., scalers), model files (\texttt{.joblib} for classical models and checkpoints for deep learning), and diagnostic figures such as PR curves, score histograms, and confusion matrices. This tracking discipline supports direct comparison across the anomaly detection and fault classification experiments.



\chapter{Anomaly Detection}

\section{Introduction}
Anomaly detection is the process of distinguishing between normal and anomalous states in a PV system. This chapter introduces the problem, 
then presents the explored models, and finally discusses the experimental results.


\chapter{Fault Diagnosis}\label{ch:fault-diagnosis}

\section{Introduction}

After an anomalous operating state is detected, the next step in a Fault Detection and Diagnosis (FDD) workflow is to identify the most likely fault type. In practical operation, fault diagnosis supports decision-making by narrowing the set of plausible causes and enabling targeted inspection and corrective actions. In this project, diagnosis is formulated as a supervised multi-class classification task performed on engineered features extracted from the PV system measurements.

Unlike the anomaly detection stage, which aims to separate healthy behavior from abnormal behavior using mostly normal-only training data (Chapter~\ref{ch:anomaly-detection}), the diagnosis stage assumes that labeled fault categories are available for training. The objective is therefore not to produce an anomaly score, but to output a discrete fault class among a predefined set of operating states.

\section{Task Formulation and Methodology}

\paragraph{Problem definition.}
Let $x \in \mathbb{R}^F$ denote a feature vector computed from a time segment of PV electrical and meteorological signals, and let $y \in \{0,\dots,C-1\}$ denote the corresponding fault label among $C$ classes. The diagnosis task consists in learning a classifier $f_\theta$ that maps $x$ to a class probability vector $\hat{p}(y\mid x)$, and predicts the most likely class $\hat{y}=\arg\max_c \hat{p}(c\mid x)$. Because fault classes can be imbalanced in realistic PV logs, the training and evaluation protocol must account for class frequency differences.

\paragraph{Inputs.}
Diagnosis is performed on pre-engineered tabular features rather than on raw time-series windows. Feature computation is executed upstream in the preprocessing and featurization pipeline, and the classification stage consumes the resulting feature tables. This design keeps the diagnosis models lightweight and enables transparent analysis of feature contributions.

\paragraph{Evaluation protocol and metrics.}
The dataset is split into training, validation, and test partitions using the same leakage-aware splitting protocol described in Section~\ref{sec:implementation-splitting}. Model selection is primarily driven by macro-averaged F1-score (macro F1), which assigns equal weight to each class and is therefore sensitive to minority-class performance. For completeness, additional metrics are reported, including overall accuracy, balanced accuracy, macro and weighted precision/recall/F1, and one-vs-rest (OvR) precision--recall AUC values (weighted and per-class).

\section{Shared Experimental Infrastructure}\label{sec:diag-shared-infra}

This section summarizes the shared training and evaluation infrastructure used by all implemented diagnosis models. The common utilities are implemented in the module \texttt{src/modeling/classification/ml/common.py}, and are designed to enforce consistent validation discipline, reproducibility, and comparable reporting across models.

\paragraph{Pipeline overview.}
The diagnosis workflow follows a fixed sequence of stages: data ingestion, split generation, preprocessing, feature extraction, and classifier training. In this chapter, only the classification stage is detailed. The classifier is trained and evaluated exclusively on the engineered feature tables produced by the upstream pipeline.

\paragraph{Feature ingestion and run management.}
Feature tables are loaded through a single entrypoint that resolves the latest (or a specified) feature-generation run for the chosen dataset and feature profile. This mechanism ensures that model runs are traceable to a concrete, versioned feature artifact, and it allows repeatable comparisons across classifiers without re-running feature extraction.

\paragraph{Threading policy.}
To make experimentation efficient and predictable across machines, a resource-aware threading policy is applied. The available CPU budget is detected automatically, one core is reserved for the operating system, and the remaining threads are split between parallel hyperparameter optimization trials and per-trial model training threads. Command-line overrides allow the user to bound the total thread usage and the number of parallel optimization jobs.

\paragraph{Label encoding.}
Fault labels are encoded into contiguous integer identifiers $\{0,\dots,C-1\}$ using a fitted label encoder. The encoder is trained on the training labels and then applied unchanged to validation and test partitions. It is serialized alongside the model artifact to ensure that inference uses the same class mapping as training.

\paragraph{Hyperparameter optimization (HPO) and validation modes.}
All implemented diagnosis models support automatic hyperparameter optimization through a unified interface. Two validation modes are available. The default mode uses a classical holdout protocol in which each trial is fitted on the training partition and scored on the validation partition using macro F1. An alternative mode performs a per-class, segment-aware temporal cross-validation in which folds follow an expanding-window structure. This second mode provides a stronger guard against temporal leakage when fault segments exhibit time-dependent structure.

\paragraph{Probability calibration.}
In addition to maximizing classification accuracy, reliable probability estimates are important for downstream decision-making (e.g., confidence-based alerting and operator interfaces). Therefore, an optional probability calibration stage is enabled by default. After a base classifier is fitted on the training partition, a calibration model is fitted on the validation partition using sigmoid calibration (Platt scaling). Both uncalibrated and calibrated probability outputs are evaluated on the test set. Calibration quality is summarized with metrics such as log loss, a multi-class Brier score, and bin-based calibration errors including the Expected Calibration Error (ECE) and Maximum Calibration Error (MCE).

\paragraph{Leakage checks.}
Before and after training, automated checks are executed to detect common leakage failure modes. In particular, an exact-duplicate overlap check verifies that identical feature rows do not appear across training and validation partitions, and a leakage report is produced on the test partition unless explicitly skipped. These controls complement the splitting protocol by providing empirical evidence that partition boundaries are respected.

\paragraph{Reporting artifacts.}
To facilitate interpretability and comparison, all models produce a consistent reporting suite. This includes a confusion matrix (raw counts and normalized), OvR precision--recall curves per class, a predicted-class distribution plot, and reliability diagnostics when calibration is enabled. For tree-based models that expose feature importance scores, the top-ranked features are also reported to support qualitative analysis of which engineered variables drive the diagnosis decisions.

\section{Model Lineup and Results}\label{sec:diag-model-lineup}

This section presents the classification models implemented for fault diagnosis. All models operate on the same engineered feature representation and are trained and evaluated under the shared infrastructure described in Section~\ref{sec:diag-shared-infra}.

\begin{table}[H]
\centering
\caption{Summary of implemented diagnosis models and their shared training protocol. HPO: Hyperparameter Optimization.}\label{tab:diag-model-summary}
\begin{tabular}{lccc}
\toprule
\textbf{Model} & \textbf{Family} & \textbf{Class balancing} & \textbf{Calibration}\\
\midrule
LightGBM      & Gradient-boosted trees & Balanced class weights & Sigmoid on validation\\
CatBoost      & Gradient-boosted trees & Balanced class weights & Sigmoid on validation\\
ExtraTrees    & Randomized tree ensemble & Balanced class weights & Sigmoid on validation\\
\bottomrule
\end{tabular}
\end{table}

\subsection{LightGBM Classifier}

LightGBM is a gradient-boosting framework that builds an ensemble of decision trees sequentially, where each new tree focuses on correcting the errors made by the current ensemble. In the diagnosis setting, the model learns non-linear decision boundaries over engineered features and can capture complex interactions between electrical indicators and contextual variables. LightGBM is selected as the default classifier because it provides a strong accuracy--efficiency trade-off on tabular data and scales well during hyperparameter search.

To mitigate class imbalance, the training procedure applies automatic class weighting so that minority fault classes receive proportionally higher penalty when misclassified. Hyperparameters such as the learning rate, number of boosting rounds, and tree complexity are tuned using the shared HPO protocol. When probability calibration is enabled, a sigmoid calibrator is fitted on validation predictions to improve the alignment between predicted confidence and empirical correctness.

\paragraph{Results}
% Intentionally left blank (to be filled with experimental results and analysis).

\subsection{CatBoost Classifier}

CatBoost is another gradient-boosting method that uses symmetric trees and an ordered boosting strategy designed to reduce prediction shift during training. Although CatBoost is particularly known for robust handling of categorical features, it also performs competitively on fully numerical feature sets and offers stable optimization behavior. In this project, CatBoost is implemented as an alternative to LightGBM within the same training harness.

The model is tuned with the same primary objective (macro F1) and the same validation modes. Class imbalance is handled through the algorithm's native balanced class-weighting mechanism. When enabled, the same sigmoid calibration stage is applied to CatBoost probability outputs using validation data.

\paragraph{Results}
% Intentionally left blank (to be filled with experimental results and analysis).

\subsection{ExtraTrees Classifier}

ExtraTrees (Extremely Randomized Trees) is an ensemble method that averages the predictions of many decision trees built with strong randomization. In contrast to boosting, trees in ExtraTrees are trained largely independently, and randomness is injected by selecting split thresholds at random. This often yields robust performance with limited tuning effort, and it provides a useful non-boosting baseline for engineered-feature classification.

In the implemented pipeline, the number of trees, depth limits, split constraints, and feature subsampling strategy are tuned under the shared HPO protocol. Class imbalance is handled via balanced class weights. Probability calibration, when enabled, follows the same validation-based sigmoid calibration procedure.

\paragraph{Results}
% Intentionally left blank (to be filled with experimental results and analysis).


\section{Conclusion}
This chapter implemented and evaluated a two-stage Fault Detection and Diagnosis (FDD) pipeline on the Costa PV dataset. For anomaly detection, \textbf{PC-Flow (baseline\_raw)} is the winning model, achieving a macro per-class PR-AUC of $0.990 \pm 0.005$ and worst-class PR-AUC of $0.962 \pm 0.021$ under 5-fold episode-stratified cross-validation. The ablation study demonstrates that physics conditioning drives the performance gain while the normalizing flow's internal representation learning makes additional hand-engineered features unnecessary. For fault diagnosis, \textbf{ExtraTrees} is the winning classifier, reaching a macro F1 of $0.959 \pm 0.078$, the highest among the three gradient-boosted and randomised tree baselines, with the lowest cross-fold variance. Together, PC-Flow and ExtraTrees form a coherent FDD pipeline: PC-Flow provides sensitive, well-ranked anomaly alarms via exact conditional density estimation, and ExtraTrees gives precise fault identification supported by interpretable feature importances that confirm the value of the handcrafted physics-informed feature profile.



